
\newcommand{\icepick}{Schr\"{o}dinger's ice-pick}

\begin{question}
  Estimate the rough order of magnitude of the length of time that an
  ice pick can be balanced on its point if the only limitation is that
  set by the Heisenberg uncertainty principle. Assume that the point
  is sharp and that the point and the surface on which it rests are
  hard. You may make approximations which do not alter the general
  order of magnitude of the result. Assume reasonable values for the
  dimensions and weight of the ice pick. Obtain an approximate
  numerical result and express it \emph{in seconds}.
\end{question}

\begin{purpose}
  This is a classical problem in high-school physics: connect the
  uncertainty relation to the real-world.  We can also see this with
  things like entanglement in macro-scale systems (I think they
  mentioned a pendulum somewhere in the book) and Bose condensates.

  There are several concepts which need to come together for this
  problem

  \begin{itemize}
  \item Application of the $X/P$ uncertainty relation.
  \item Application of the $E/t$ uncertainty relation.
  \item Taking expectation values using wave functions.
  \item The Wigner function.
  \item Wave function collapse after observations.
  \item Mass energy and momentum implications from special relativity.
  \end{itemize}

  Ultimately, this problem ties together just about all of the
  important topics in chapter 1 and then some.  We're going to solve
  the problem using a few different approaches:
  
  \begin{questionpart}
    The most straight-forward normal way of doing this problem is to
    assume a wave function for the position, determine the \acp{pdf}
    for $P$ and $X$ and integrate against the fall-time as a function
    of those variables.  True to form, things are out of order here so
    we'll have to skip around a bit, starting with working problem
    \ref{1.36:theQuestion}.
  \end{questionpart}

  \begin{questionpart}
    In our first solution, we essentially use quantum mechanics to set
    up the probability distributions for the initial conditions, then
    we use classical mechanics for time-propagation.  Here, we'll use
    quantum mechanics to propagate our gaussian wave packet.
  \end{questionpart}

  \begin{questionpart}
    The next fun way to solve the problem is to leverage the $E/t$
    uncertainty relation to determine the amount of rotational kinetic
    energy in the system.  Using the time expended during the
    observation of energy, we then reverse-engineer the presumed
    initial conditions and compute the expecation value just as we did
    in the previous part.
  \end{questionpart}

  Note that I actually solved this problem out of order.  Problem
  \ref{1.36:theQuestion} involves proving one of the key relations for
  the solution to this problem so I decided to solve that one first.

\end{purpose}

\begin{answer}
  So, if we had a perfect and perfectly-determined world, then the
  ice-pick would be perfectly vertical, made of Hyle, infinitely
  rigid, no mechanical/electrical/use-the-force-luke perturbations and
  it would just sit there until the heat-death of the universe.
  Newsflash to anyone watching Krasnov and global warming...we don't
  live in a perfect world.  The fundamental indeterminability of
  quantum mechanics ensures that we are working with \icepick{}.  That
  means that whatever indeterminability there is results in the
  ice-pick being in a superposition of all of those states.

  If everything is in a superposition of states then it doesn't really
  make sense to ask how long it will take for that specific trial
  since it is unknowable.  What we \emph{can} ask, is what the
  expectation-value is for the fall-time.  But even that statement
  raises questions.  If our infinite set of monkeys with their
  infinite typewriters repeated the experiment an infinite number of
  times would one of them take an infinite amount of time to fall?
  While we're on that topic, what does the ``fall-time'' even mean?

  Taking the last part first and the first part last, ``fall-time'' is
  defined as when a recording-device of some kind recognizes a signal
  which injects all sorts of statistical issues.  When are there
  enough electrons on the gate of the Field Effect Transistors in the
  CMOS logic sufficiently high to get past the ohmic region?  But we
  understand that statistical analysis reasonably-well -- it's a bunch
  of gaussians with $\sigma$'s that we have determined experimentally
  (assuming it happens in time-ranges $\ge 10^{-10}$s).  The more
  interesting quantum mechanics question is about what it means for
  \icepick{} to trigger the sensor.  When \icepick{} falls, its
  gaussian wave-packet is moving and expanding.  Initial
  indeterminability of the momentum (becoming $\omega_0$) and position
  (becoming $\theta_0$) leads to time-expanding indeterminability of
  the position as it falls.  So we can define our ``fall-time'' as the
  time at which the expectation-value of the spatial wave funtion
  (reparameterized for $\theta$) reaches the ground (i.e. $\theta =
  \pi/2$).  For simplicity, though, since time-evolution of
  wave-functions isn't something that we covered, we're going to use
  quantum mechanics to determine initial conditions and then (FIXME:
  lagrangian or hamiltonian?) classical mechanics to determine the
  fall-time from those initial conditions.

  Finally, the infinite monkeys...well..they eventually \emph{do}
  cause \icepick{} to balance for an arbitrarily-long time.


  \begin{answerpart}{The ``normal'' way}

    I mean...I \emph{assume} this is the normal way at least.  Like
    the old joke about Werner Heisenberg getting pulled over by the
    cops goes, we can't know our position and momentum at the same
    time.  We just went to great pains to prove that if $\vec{a}^{*}
    \cdot \vec{a} \ge 0$ then the uncertainty relation is a thing, so
    let's use that as our starting point.

    We're going to need to pick a wave function to work with and the
    \ac{gwp} fits the bill for a couple of reasons:

    \begin{enumerate}
    \item They are known to minimize the $\hat{X}/\hat{P}$ uncertainty
      (i.e. $\dispersion{\hat{X}}\dispersion{\hat{P}} = \frac{1}{4}
      \abs{\expval{[\hat{X}, \hat{P}]}}^2$

    \item The problem setup actually kinda clues us into the concept
      of using a Gaussian in that it seems quite natural to assume
      that any machine which is capable of placing \icepick{} is
      probably going to have a Gaussian in its placement.
    \end{enumerate}

    Since this is the first problem in which we need to really dig
    into this wave function, we're going to show the full derivation
    below then record the useful results in our \ac{qrf} section
    \ref{q:wavey:gwp}.
    
    \begin{equation}
      \tag{1.146}
      \dispersion{A}\dispersion{B} \ge \frac{1}{4} \abs{\expval{[A, B]}}^2
    \end{equation}

    So swapping out those operators to be the ones we're interested in:
  
    \begin{equation}
      \begin{split}
        \dispersion{x}\dispersion{p}
        &\ge \frac{1}{4} \abs{\expval{[x, p]}}^2\\
        &= \frac{1}{4} \abs{\expval{i\hbar}}^2\\
        &= \frac{\hbar^2}{4}\\
      \end{split}
    \end{equation}

    So, if we have equality, then it means that the only wave
    functions we're going to accept are those which meet that
    equality.  Here, we're going to make an assumption.  Since the
    whole concept is ``place the ice pick and make it reeeeeeally
    steady'' then we're going to assume that placement of the ice-pick
    is a gaussian with respect to the top of the ice pick.  It's also
    worth noting that Gaussian wave packets are known to minimize
    uncertainty.

    \begin{equation}
      \tag{1.267}
      \braket{x'}{\alpha}
      = \left[\frac{1}{\pi^{1/4}\sqrt{d}}\right]
      \exp[ikx' - \frac{x'^{2}}{2d^2}]
    \end{equation}

    With our spatial wave-function in hand, we now have 4 values to
    compute: $\expval{\hat{X}}$, $\expval{\hat{X}^2}$,
    $\expval{\hat{P}}$, and $\expval{\hat{P}^2}$.  Let's start with
    $\expval{\hat{X}}$.  We're going to use $\psi$ to represent the
    spatial wave-function rather than the more general $\alpha$.

    \begin{equation}
      \begin{split}
        \expval{\hat{X}}
        &= \mel{\psi}{\hat{X}}{\psi}\\
        &= \mel{\psi}{\hat{X}\hat{1}}{\psi}\\
        &= \int dx' (\bra{\psi})(\hat{X}\dyad{x'})(\ket{\psi})\\
        &= \int dx' \mel{\psi}{\hat{X}}{x'}\braket{x'}{\psi}\\
        &= \int dx' x'\braket{\psi}{x'}\braket{x'}{\psi}\\
        &= \int dx' x'\braket{x'}{\psi}^{*}\braket{x'}{\psi}\\
      \end{split}
    \end{equation}

    Bringing in our definition of our spatial wave funtion, we find
    that
    
    \begin{equation}
      \inlabel{exp:x1}
      \begin{split}
        \expval{\hat{X}}
        &= \int dx' x'\braket{x'}{\psi}^{*}\braket{x'}{\psi}\\
        &= \left[\frac{1}{\pi^{1/4}\sqrt{d}}\right]^2
        \int dx' x'
        \exp[\cancel{-ikx'} - \frac{x'^{2}}{2d^2}]
        \exp[\cancel{ikx'} - \frac{x'^{2}}{2d^2}]\\
        &= \frac{1}{\sqrt{\pi}d}
        \int dx' x'
        \exp[- \cancel{2}\frac{x'^{2}}{\cancel{2}d^2}]\\
        &= \frac{1}{\sqrt{\pi}d}
        \int dx' x'
        \exp[- \frac{x'^{2}}{d^2}]\\
        &= 0\\
      \end{split}
    \end{equation}

    Using equation \ref{q:integral:xgaussian}, we see that the
    expectation value for position is $0$.  So, this is a long-winded
    way of saying something we already knew to be true: the
    expectation value of a symmetric function centerd about $0$
    is... well... $0$.  You can also shortcut the integral by noting
    that you're integrating an odd function with symmetrical bounds
    which is always $0$ ($\int_{-a}^a dx x \cdot \text{symmetric}(x) =
    0$).  Let's next look at the expectation value of the square.
    We're only going to get a second $x'$ in that integral but let's
    do it anyway for completeness.

    \begin{equation}
      \begin{split}
        \expval{\hat{X}^2}
        &= \expval{\hat{X}\hat{X}}\\
        &= \mel{\psi}{\hat{X}\hat{X}}{\psi}\\
        &= \mel{\psi}{\hat{X}\hat{1}\hat{X}}{\psi}\\
        &= \int dx' (\bra{\psi})(\hat{X}\dyad{x'}\hat{X})(\ket{\psi})\\
        &= \int dx' \mel{\psi}{\hat{X}}{x'}\mel{x'}{\hat{X}}{\psi}\\
        &= \int dx' x'^2\braket{\psi}{x'}\braket{x'}{\psi}\\
        &= \int dx' x'^2\braket{x'}{\psi}^{*}\braket{x'}{\psi}\\
      \end{split}
    \end{equation}

    Once again employing integral-by-table and referencing equation
    \ref{q:integral:xxgaussian}, we find that

    \begin{equation}
      \inlabel{exp:x2}
      \begin{alignedat}{2}
        & &
        \int_{-\infty}^{\infty} dx\ x^2e^{-\alpha x^2}
        &= \frac{1}{2\alpha}\sqrt{\frac{\pi}{\alpha}}\\
        \implies & &
        \frac{1}{\sqrt{\pi}d}
        \int_{-\infty}^{\infty} dx\ x^2\exp(-\frac{x^2}{d^2})
        &= \cancel{\frac{1}{\sqrt{\pi}d}}\frac{d^2}{2}\cancel{\sqrt{d^2\pi}}\\
        &&&= \frac{d^2}{2}\\
      \end{alignedat}
    \end{equation}

    Next up, let's compute the expectation value for the momentum.
    
    \begin{equation}
      \begin{split}
        \expval{\hat{P}}
        &= \mel{\psi}{\hat{P}}{\psi}\\
        &= \mel{\psi}{\hat{1}\hat{P}}{\psi}\\
        &= \int dx' (\bra{\psi})(\dyad{x'}\hat{P})(\ket{\psi})\\
        &= \int dx' \braket{\psi}{x'}\mel{x'}{\hat{P}}{\psi}\\
      \end{split}
    \end{equation}

    Pulling in an identity from the book

    \begin{equation}
      \tag{1.249}
      \mel{x'}{\hat{p}}{\alpha} = -i \hbar \pdv{x'} \braket{x'}{\alpha}
    \end{equation}

    \begin{equation}
      \begin{split}
        \expval{\hat{P}}
        &= \int dx'
        \braket{\psi}{x'}
        \mel{x'}{\hat{P}}{\psi}\\
        &= \int dx'
        \braket{x'}{\psi}^{*}
        \left(-i \hbar \pdv{x'}\right)
        \braket{x'}{\psi}\\
        &= \left[\frac{1}{\pi^{1/4}\sqrt{d}}\right]^2
        \int dx'
        \exp[-ikx' - \frac{x'^{2}}{2d^2}]
        \left(-i \hbar \pdv{x'}\right)
        \exp[ikx' - \frac{x'^{2}}{2d^2}]\\
        &= \frac{1}{d\sqrt{\pi}}
        \int dx'
        \exp[-ikx' - \frac{x'^{2}}{2d^2}]
        (-i \hbar)
        \left(
        ik - \frac{\cancel{2}x'}{\cancel{2}d^2}
        \right)
        \exp[ikx' - \frac{x'^{2}}{2d^2}]\\
        &= \frac{\hbar}{d\sqrt{\pi}}
        \int dx'
        \exp[-ikx' - \frac{x'^{2}}{2d^2}]
        \left(
        k + i\frac{x'}{d^2}
        \right)
        \exp[ikx' - \frac{x'^{2}}{2d^2}]\\
        &= \frac{\hbar}{d\sqrt{\pi}}
        \int dx'
        \left(k + i\frac{x'}{d^2}\right)
        \exp[- \frac{x'^{2}}{d^2}]\\
        &= \frac{\hbar}{d\sqrt{\pi}}
        \left\{
        k\int dx'
        \exp[- \frac{x'^{2}}{d^2}]
        + i\int dx'
        \frac{x'}{d^2}\exp[- \frac{x'^{2}}{d^2}]
        \right\}\\
      \end{split}
    \end{equation}

    Once again drawing upon our integrals in the \ac{qrf},
    specifically equations \ref{q:integral:xgaussian} and
    \ref{q:integral:xxgaussian}:

    \begin{equation}
      \inlabel{exp:p1}
      \begin{split}
        \expval{\hat{P}}
        &= \int dx'
        \braket{\psi}{x'}
        \mel{x'}{\hat{P}}{\psi}
        \\
        &= \frac{\hbar}{d\sqrt{\pi}}
        \left\{
        k\int dx'
        \exp[- \frac{x'^{2}}{d^2}]
        + i\int dx'
        \frac{x'}{d^2}\exp[- \frac{x'^{2}}{d^2}]
        \right\}
        \\
        &= \frac{k\hbar}{d\sqrt{\pi}}
        \int dx'
        \exp[- \frac{x'^{2}}{d^2}]
        \\
        &= \frac{k\hbar}{d\sqrt{\pi}}
        \sqrt{d^2\pi}
        \\
        &= k\hbar
        \\
      \end{split}
    \end{equation}
    
    So, that works out quite conveniently doesn't it.  You'll note
    that the $k$ is the imaginary term in the exponential which
    represents the momentum term.  I'm preeeettty sure the gaussian
    wave packet formulation was specifically designed to separate out
    the spatial and momentum contributions.  It also makes it pretty
    easy to see that this wave function has both a spatial and a
    momentum component and thus represents a complete state
    description.  One expectation value left.

    \begin{equation}
      \begin{split}
        \expval{\hat{P}^2}
        &= \mel{\psi}{\hat{P}\hat{P}}{\psi}\\
        &= \mel{\psi}{\hat{P}\hat{1}\hat{P}}{\psi}\\
        &= \int dx' (\bra{\psi})(\hat{P}\dyad{x'}\hat{P})(\ket{\psi})\\
        &= \int dx' \mel{\psi}{\hat{P}}{x'}\mel{x'}{\hat{P}}{\psi}\\
        &= \int dx' \mel{x'}{\hat{P}}{\psi}^{*}\mel{x'}{\hat{P}}{\psi}\\
        &=
        \int dx'
        \left(i \hbar \pdv{x'}\right)
        \braket{x'}{\psi}^{*}
        \left(-i \hbar \pdv{x'}\right)
        \braket{x'}{\psi}
        \\
        &=
        \left[\frac{1}{\pi^{1/4}\sqrt{d}}\right]^2
        \int dx'
        \left\{
        \left(i \hbar \pdv{x'}\right)
        \exp[-ikx' - \frac{x'^{2}}{2d^2}]
        \right\}
        \left\{
        \left(-i \hbar \pdv{x'}\right)
        \exp[ikx' - \frac{x'^{2}}{2d^2}]        
        \right\}
        \\
      \end{split}
    \end{equation}

    Leveraging some intermediate products from before (and being wary
    of sign errors) we find that:

    \begin{equation}
      \begin{split}
        \expval{\hat{P}^2}
        &=
        \left[\frac{1}{\pi^{1/4}\sqrt{d}}\right]^2
        \int dx'
        \left\{
        \left(i \hbar \pdv{x'}\right)
        \exp[-ikx' - \frac{x'^{2}}{2d^2}]
        \right\}
        \left\{
        \left(-i \hbar \pdv{x'}\right)
        \exp[ikx' - \frac{x'^{2}}{2d^2}]        
        \right\}
        \\
        &=
        \frac{1}{d\sqrt{\pi}}
        \int dx'
        (i \hbar)
        \left(
        -ik - \frac{\cancel{2}x'}{\cancel{2}d^2}
        \right)
        \exp[-ikx' - \frac{x'^{2}}{2d^2}]
        (-i \hbar)
        \left(
        ik - \frac{\cancel{2}x'}{\cancel{2}d^2}
        \right)
        \exp[ikx' - \frac{x'^{2}}{2d^2}]
        \\
        &=
        \frac{\hbar^2}{d\sqrt{\pi}}
        \int dx'
        \left(
        k - i\frac{x'}{d^2}
        \right)
        \left(
        k + i\frac{x'}{d^2}
        \right)
        \exp[-\frac{x'^{2}}{d^2}]
        \\
        &=
        \frac{\hbar^2}{d\sqrt{\pi}}
        \int dx'
        \left(
        k^2 + \frac{x'^2}{d^4}
        \right)
        \exp[-\frac{x'^{2}}{d^2}]
        \\
        &=
        \frac{\hbar^2}{d\sqrt{\pi}}
        \left[
        k^2
        \int dx'
        \exp[-\frac{x'^{2}}{d^2}]
        +
        \int dx'
        \frac{x'^2}{d^4}
        \exp[-\frac{x'^{2}}{d^2}]
        \right]
        \\
      \end{split}
    \end{equation}

    Once again bringing out our table of integrals

    \begin{equation}
      \inlabel{exp:p2}
      \begin{split}
        \expval{\hat{P}^2}
        &=
        \frac{\hbar^2}{d\sqrt{\pi}}
        \left[
        k^2
        \int dx'
        \exp[-\frac{x'^{2}}{d^2}]
        +
        \int dx'
        \frac{x'^2}{d^4}
        \exp[-\frac{x'^{2}}{d^2}]
        \right]
        \\
        &=
        \frac{\hbar^2}{d\sqrt{\pi}}
        \left[
        k^2
        d\sqrt{\pi}
        +
        \frac{1}{d^4}
        \frac{d^2}{2}d\sqrt{\pi}
        \right]
        \\
        &=
        \frac{\hbar^2}{\cancel{d\sqrt{\pi}}}
        \cancel{d\sqrt{\pi}}
        \left(
        k^2
        +
        \frac{1}{2d^2}
        \right)
        \\
        &=
        \hbar^2
        \left(
        k^2
        +
        \frac{1}{2d^2}
        \right)
        \\
      \end{split}
    \end{equation}

    \begin{equation}
      \inlabel{exp:combined}
      \begin{alignedat}{2}
        \expval{X}^2 &= 0 & \quad{} &\text{(eq \inref{exp:x1})}\\
        \expval{X^2} &= \frac{d^2}{2} & &\text{(eq \inref{exp:x2})}\\
        \expval{P}^2 &= \hbar^2k^2 & &\text{(eq \inref{exp:p1})}\\
        \expval{P^2} &=
        \hbar^2
        \left(
        k^2
        +
        \frac{1}{2d^2}
        \right)
        & &\text{(eq \inref{exp:p2})}\\
      \end{alignedat}
    \end{equation}
    
    recalling our definition of the variance

    \begin{equation*}
      \tag{1.144}
      \dispersion{A} = \expval{A^2} - \expval{A}^2
    \end{equation*}

    \begin{equation}
      \begin{split}
        \dispersion{X}\dispersion{P}
        &=
        \left(\expval{\hat{X}^2} - \expval{\hat{X}}^2\right)
        \left(\expval{\hat{P}^2} - \expval{\hat{P}}^2\right)
        \\
        &=
        \left(\frac{d^2}{2} - 0\right)
        \left(
        \hbar^2
        \left(
        k^2
        +
        \frac{1}{2d^2}
        \right)
        - \hbar^2k^2
        \right)
        \\
        &=
        \hbar^2
        \frac{\cancel{d^2}}{2}
        \left(
        \cancel{k^2}
        +
        \frac{1}{2\cancel{d^2}}
        - \cancel{k^2}
        \right)
        \\
        &= \frac{\hbar^2}{4}
      \end{split}
    \end{equation}

    Cool...  Well... I guess we just proved that the gaussian wave
    packet does, indeedily-do, minimize the uncertainty.  So, what
    does this tell us?  Well, for starters, it tells us that we made a
    good selection of wave function for our experiment.  Next, it
    gives us some equations to work with that have two variables ($d$
    and $k$) with one representing position indeterminability and the
    other representing group momentum.  As we see above, the group
    velocity doesn't impact the overall indeterminability as it
    cancels out (of course, so does the position indeterminability).
    The $e^{ikx'}$ term is clearly recognizable as both a phasor and
    as an independent operator.  We actually see this operator later
    in problem \ref{1.36:theQuestion} where we call it the Jerk
    Operator.

    This is where we make another simplifying assumption: set $k = 0$.
    Since we're trying our best to keep \icepick{} standing, we're
    going to assume that any momentum is entirely rotational in
    nature.  As such, we can safely set the group momentum to be 0.
    Substituting those values into our expectations from before, we
    find that

    \begin{equation}
      \inlabel{exp:simplified}
      \begin{alignedat}{2}
        \expval{X}^2 &= 0 & \quad{} &\text{(eq \inref{exp:x1})}\\
        \expval{X^2} &= \frac{d^2}{2} & &\text{(eq \inref{exp:x2})}\\
        \expval{P}^2 &= 0 & &\text{(eq \inref{exp:p1} modified)}\\
        \expval{P^2} &= \frac{\hbar^2}{2d^2} & &\text{(eq \inref{exp:p2} modified)}\\
      \end{alignedat}
    \end{equation}
    
    This is where we run into an issue with the concept of an
    expectation-value which is likely to require some limit trickery.
    The denominator in the $\ln$ includes both $\theta_0$ and
    $\omega_0$ which means that it blows up as you approach 0 on both
    of them.  There is one thing that, physically, might help us out
    here -- $P(\omega_0)|_{x_0 \approx 0} \approx 0$.  So let's see if
    that integral converges.  To find out, we'll start by computing
    the probabilities of $x'$ and $p'$.  \S 1.7.3 in the book
    discusses this exact issue and formulates the position and
    momentum space wave functions (which happen to be fourier
    transforms of one another... iiinteresting... kinda makes you
    wonder about Shor's algorithm doesn't it).

    Starting with the statistical interpretation from the text just
    after eq 1.256, ``the probability that a measurement of $p$ gives
    eigenvalue $p'$ within a narrow interval $dp'$ is
    $\abs{\braket{p'}{\alpha}}^2dp'$.  It is customary to call
    $\braket{p'}{\alpha}$ the \textbf{momentum-space wave function};
    the notation $\phi_{\alpha}(p')$ is often used''.  So, if we have
    an expression for $\phi_{\alpha}(p')$ then we can find the
    probability:

    \begin{equation}
      P(p = p' \pm dp'/2) = \abs{\phi_{\alpha}(p')}^2dp'
    \end{equation}

    Later, Sakurai provides us with that exact equation:
    
    \begin{equation}
      \tag{1.266a}
      \psi_{\alpha}(x')
      = \left[ \frac{1}{\sqrt{2\pi\hbar}} \right]
      \int dp' \exp(\frac{ip'x'}{\hbar})\phi_{\alpha}(p')
    \end{equation}

    \begin{equation}
      \tag{1.266b}
      \phi_{\alpha}(p')
      = \left[ \frac{1}{\sqrt{2\pi\hbar}} \right]
      \int dx' \exp(\frac{-ip'x'}{\hbar})\psi_{\alpha}(x')
    \end{equation}

    Let's use those equations to get ourselves a referenceable
    expression for $\phi(p')$ and $\psi(x')$.  $\psi$ is, of course,
    pretty easy since we started with an ansatz'd sapatial
    wave-function deciding to infer the momentum-space wave-function
    from there.  Going back to 1.267 and re-stating here with our own
    internal numbering that we can more easily reference but also
    inserting our assumption that $k=0$ we find that

    \begin{equation}
      \inlabel{psi}
      \begin{split}
        \psi(x')
        &= \braket{x'}{\alpha}\\
        &= \left[\frac{1}{\pi^{1/4}\sqrt{d}}\right]
        \exp[-\frac{x'^{2}}{2d^2}]\\
      \end{split}
    \end{equation}
    
    Evaluating $\phi$ we have

    \begin{equation}
      \begin{split}
        \phi(p')
        &= 
        \left[ \frac{1}{\sqrt{2\pi\hbar}} \right]
        \int dx' \exp(\frac{-ip'x'}{\hbar})\psi(x')
        \\
        &= 
        \left[ \frac{1}{\sqrt{2\pi\hbar}} \right]
        \left[\frac{1}{\pi^{1/4}\sqrt{d}}\right]
        \int dx'
        \exp(\frac{-ip'x'}{\hbar})
        \exp[-\frac{x'^{2}}{2d^2}]
        \\
        &= 
        \left[ \frac{1}{\sqrt{2\pi\hbar}} \right]
        \left[\frac{1}{\pi^{1/4}\sqrt{d}}\right]
        \int dx'
        \exp(\frac{-ip'x'}{\hbar}
        - \frac{x'^{2}}{2d^2})
        \\
      \end{split}
    \end{equation}

    Pulling in equation 3.323.2 from our integral table.
    \cite[eq3.323.2]{gr}

    \begin{equation}
      \tag{3.323.2}
      \int_{-\infty}^{\infty} \exp\left(-p^2x^2 \pm qx\right) dx
      = \exp\left(\frac{q^2}{4p^2}\right)\frac{\sqrt{\pi}}{p}
    \end{equation}

    Substituting into our expression above we find that we set $p =
    \frac{1}{d\sqrt{2}}$ and $q = -ip'/\hbar$ to find that


    \begin{equation}
      \inlabel{phi}
      \begin{split}
        \phi(p')
        &= 
        \left[ \frac{1}{\sqrt{2\pi\hbar}} \right]
        \int dx' \exp(\frac{-ip'x'}{\hbar})\psi(x')
        \\
        &= 
        \left[ \frac{1}{\sqrt{2\pi\hbar}} \right]
        \left[\frac{1}{\pi^{1/4}\sqrt{d}}\right]
        \int dx'
        \exp(\frac{-ip'x'}{\hbar}
        - \frac{x'^{2}}{2d^2})
        \\
        &= 
        \left[ \frac{1}{\sqrt{2\pi\hbar}} \right]
        \left[\frac{1}{\pi^{1/4}\sqrt{d}}\right]
        \exp(\frac{q^2}{4p^2})
        \frac{\sqrt{\pi}}{p}
        \\
        &= 
        \left[ \frac{1}{\sqrt{\cancel{2\pi}\hbar}} \right]
        \left[\frac{1}{\pi^{1/4}\sqrt{d}}\right]
        \exp(-\frac{p'^2}{4\hbar^2}2d^2)
        d\cancel{\sqrt{2\pi}}
        \\
        &= 
        \frac{1}{\pi^{1/4}}
        \sqrt{\frac{d}{\hbar}}
        \exp[-\frac{1}{2}\left(\frac{p'd}{\hbar}\right)^2]
        \\
        &= 
        \frac{1}{\pi^{1/4}}
        \sqrt{\frac{d}{\hbar}}
        \exp[-\frac{d^2}{2\hbar^2}p^2]
        \\
      \end{split}
    \end{equation}

    Now that we have expressions for $\psi(x')$ and $\phi(p')$ it's
    worth taking a moment to consider the gaussian nature of them.  A
    ``Normal'' or ``Gaussian'' distribution in classical probability
    theory is expressed as

    \begin{equation}
      N(\mu, \sigma^2)
      = \frac{1}{\sigma\sqrt{2\pi}}\exp[-\frac{(x-\mu)^2}{2\sigma^2}]
    \end{equation}\incite[inside cover]{hogtan}

    Mapping our equations above, we find that

    \begin{equation}
      \begin{split}
        P(x')
        &= \abs{\psi(x')}^2
        = \frac{1}{d\sqrt{\pi}} \exp[-\frac{x'^{2}}{d^2}]
        = N(\mu_\psi, \sigma_\psi)
        \\
        \mu_\psi &= 0\\
        \sigma_\psi &= \frac{d}{\sqrt{2}}\\
        \\
        P(p')
        &= \abs{\phi}^2
        = \frac{1}{\sqrt{\pi}} \frac{d}{\hbar} \exp[-\frac{d^2}{\hbar^2}p^2]
        = N(\mu_\phi, \sigma_\phi)
        \\
        \mu_\phi &= 0\\
        \sigma_\phi &= \frac{\hbar}{d\sqrt{2}}\\
      \end{split}
    \end{equation}

    So, with that parameterization, we see pretty quickly that we have
    an inverse relationship between $\sigma_\psi$ and $\sigma_\phi$
    \ldots want to make one small, ya gotta make the other big.  Armed
    with these expressions, it's time to set up our classical integral
    for the expectation value and see what happens.  Are you as
    excited as I am right now?  No?  Well if you don't like it then
    you can take a hike mister!  Ahheem... moving on.  In a classical
    expectation-value problem for the fall time ($t_f$), we would set
    it up as follows

    \begin{equation}
      \expval{t_f} = \iint_{-\infty}^{\infty} P(x)P(p)t(x, p) dx dp
    \end{equation}

    Here, we have parameterized the fall-time as a function $t(x, p)$
    where the $x$ and $p$ are really $x_0$ and $p_0$ respectively.  We
    already have expressions above for the probabilities so that just
    leaves the classical-mechanics computation of the fall-time as a
    function of initial conditions.  So...that leaves computing the
    fall-time as a function of initial conditions using classical
    mechanics... well... it turns out to be \emph{quite difficult}.
    The thing is, I'm here to do two things... solve quantum mechanics
    problems and chew bubblegum...and I'm all out of minimum-action
    bubblegum.

    I've seen other solutions to this problem which attempt to solve
    Close Encounters of the Third Kind... errm... Elliptic Integrals
    of the First kind with various
    approximations.\incite{tigersPendulum} Instead of spending another
    12 days and 12 pages to solve this on my own, I'm going to just
    let my handy-dandy AI tutor sketch out the class-mech portion for
    me.  There are going to be a few approximations we take which will
    limit the accuracy of our expectation value and I'll do my best to
    evaluate the order of magnitude of the ensuing residual.

%%     We model the ice pick as a point mass $m$ at the end of a massless
%%     rod of length $L$. Newton's second law for rotation is 
%%     $\tau = I \ddot{\theta}$. The moment of inertia is $I = mL^2$, and
%%     the torque from gravity is $\tau = mgL \sin\theta$. This gives the 
%%     exact equation of motion:

%%     \begin{equation}
%%       mL^2 \ddot{\theta} = mgL \sin\theta 
%%       \implies \ddot{\theta} - \frac{g}{L} \sin\theta = 0
%%     \end{equation}

%%     Because our quantum Gaussian wave packets aggressively constrain the
%%     initial position to microscopic perturbations, we can safely invoke 
%%     the small-angle approximation $\sin\theta \approx \theta$. Defining 
%%     $\gamma \equiv \sqrt{g/L}$, we get a linear second-order ODE:

%%     \begin{equation}
%%       \ddot{\theta} - \gamma^2 \theta = 0
%%     \end{equation}

%%     The general solution to this differential equation is a superposition 
%%     of a growing and a decaying exponential:

%%     \begin{equation}
%%       \theta(t) = C_1 e^{\gamma t} + C_2 e^{-\gamma t}
%%     \end{equation}

%%     To find the constants, we apply our initial conditions at $t=0$, where
%%     $\theta(0) = \theta_0$ and $\dot{\theta}(0) = \omega_0$:

%%     \begin{equation}
%%       \begin{alignedat}{2}
%%         \theta_0 &= C_1 + C_2 \\
%%         \omega_0 &= \gamma C_1 - \gamma C_2 \implies \frac{\omega_0}{\gamma} = C_1 - C_2
%%       \end{alignedat}
%%     \end{equation}

%%     Adding the two equations together isolates $C_1$:

%%     \begin{equation}
%%       C_1 = \frac{1}{2}\left( \theta_0 + \frac{\omega_0}{\gamma} \right)
%%     \end{equation}

%%     For any macroscopic fall time, the decaying exponential term 
%%     $e^{-\gamma t}$ vanishes almost instantly. The fall is completely
%%     dominated by the growing exponential: $\theta(t) \approx C_1 e^{\gamma t}$.
    
%%     The ice pick hits the table when the angle reaches $\pi/2$. However,
%%     if the initial momentum pushes the pick strongly toward the vertical axis,
%%     $C_1$ can become negative, meaning the pick swings through the axis 
%%     and falls to $-\pi/2$. We handle both cases gracefully by taking the 
%%     absolute value:

%%     \begin{equation}
%%       |C_1| e^{\gamma t_f} = \frac{\pi}{2}
%%     \end{equation}

%%     Solving for the fall time $t_f$:

%%     \begin{equation}
%%       \begin{split}
%%         t_f 
%%         &= \frac{1}{\gamma} \ln\left( \frac{\pi}{2|C_1|} \right) \\
%%         &= \sqrt{\frac{L}{g}} \ln\left( \frac{\pi}{\left| \theta_0 + \omega_0\sqrt{\frac{L}{g}} \right|} \right)
%%       \end{split}
%%     \end{equation}

%%     Finally, we map our initial angle and angular velocity to our quantum 
%%     phase-space variables $x'$ (arc length) and $p'$ (linear momentum). 
%%     Using $\theta_0 = x'/L$ and $\omega_0 = p'/(mL)$, we find:

%%     \begin{equation}
%%       t_f(x', p') = \sqrt{\frac{L}{g}} \ln\left( \frac{\pi}{\left| \frac{x'}{L} + \frac{p'}{mL}\sqrt{\frac{L}{g}} \right|} \right)
%%     \end{equation}





    
    

  %%   \begin{equation}
  %%     \begin{split}
  %%       \expval{t_f}
  %%       &= \iint_{-\infty}^{\infty} P(x_0=x') P(p_0=p') t_f(x', p') dx' dp'
  %%       \\
  %%       &= \iint_{-\infty}^{\infty}
  %%       \frac{1}{d\sqrt{\pi}} \exp[-\frac{x'^{2}}{d^2}]
  %%       \frac{1}{\sqrt{\pi}} \frac{d}{\hbar} \exp[-\frac{d^2}{\hbar^2}p'^2]
  %%       \sqrt{\frac{L}{g}}\ln( \frac{\pi}{\theta_0 + \omega_0 \sqrt{\frac{L}{g}}})
  %%       dx' dp'
  %%       \\
  %%     \end{split}
  %%   \end{equation}

  %%   Noting that the fall-time is parameterized for $\theta_0$ and
  %%   $\omega_0$ instead of $x'$ and $p'$, we first show that
  %%   relationship.  For the moment, we're only examining situations in
  %%   which the samll-angle approximation works.

  %%   \begin{equation}
  %%     \begin{split}
  %%       \omega_0 &\approx \frac{p_0}{mL} \\
  %%       \\
  %%       \theta_0 &\approx \frac{x_0}{L} \\
  %%     \end{split}
  %%   \end{equation}

  %%   \begin{equation}
  %%     \begin{split}
  %%       \expval{t_f}
  %%       &= \iint_{-\infty}^{\infty}
  %%       \frac{1}{d\sqrt{\pi}} \exp[-\frac{x'^{2}}{d^2}]
  %%       \frac{1}{\sqrt{\pi}} \frac{d}{\hbar} \exp[-\frac{d^2}{\hbar^2}p'^2]
  %%       \sqrt{\frac{L}{g}}\ln( \frac{\pi}{\theta_0 + \omega_0 \sqrt{\frac{L}{g}}})
  %%       dx' dp'
  %%       \\
  %%       &=
  %%       \frac{1}{\pi\hbar}
  %%       \sqrt{\frac{L}{g}}
  %%       \iint_{-\infty}^{\infty}
  %%       \exp[-\frac{x'^{2}}{d^2}]
  %%       \exp[-\frac{d^2}{\hbar^2}p'^2]
  %%       \ln( \frac{\pi}{\frac{x'}{L} + \frac{p'}{mL} \sqrt{\frac{L}{g}}})
  %%       dx' dp'
  %%       \\
  %%       &=
  %%       \frac{1}{\pi\hbar}
  %%       \sqrt{\frac{L}{g}}
  %%       \iint_{-\infty}^{\infty}
  %%       \exp[-\frac{x'^{2}}{d^2} - \frac{d^2}{\hbar^2}p'^2]
  %%       \ln( \frac{\pi}{x'\frac{1}{L} + p' \sqrt{\frac{1}{m^2gL}}})
  %%       dx' dp'
  %%       \\
  %%       &=
  %%       A
  %%       \iint_{-\infty}^{\infty}
  %%       \exp[-X^2 - P^2]
  %%       \ln( \frac{\pi}{Bx' + Cp'} )
  %%       dx' dp'
  %%       \\
  %%     \end{split}
  %%   \end{equation}
  %% \end{answerpart}

    %% I'm here to do two things... solve quantum problems and chew
    %% bubblegum...and I'm all out of bubblegum.  What I am \emph{not}
    %% here to do is explore classmech.  It has been quite a while and I
    %% really want to get through this chapter.  Straight-up, I asked
    %% Google's AI to solve this one for me.  I'll validate it later when
    %% I go back through Goldstein (that may be a while, so if someone
    %% wants to contribute then check out the appendix).  This problem
    %% has been solved many times and LLMs tend to do a decent job of
    %% regurgitating solutions.

    %% \begin{equation}
    %%   t_f =
    %%   \sqrt{\frac{L}{g}}
    %%   \ln( \frac{\pi}{\theta_0 + \omega_0 \sqrt{\frac{L}{g}}})
    %% \end{equation}

    %% We're going to redefine a variable quickly to make it a bit easier

    %% \begin{equation}
    %%   \text{let }\sqrt{\frac{L}{g}} \equiv k
    %% \end{equation}

    %% \begin{equation}
    %%   t_f = k \ln( \frac{\pi}{\theta_0 + \omega_0k})
    %% \end{equation}

    
  %% We're going to make one change here and we're going to take an
  %% absolute value.  If the momentum is towards the normal then it's
  %% going to lose momentum on the way up, and re-capture it on the way
  %% down.  In both cases, it ends up barrelling towards the ground with
  %% the original $\omega_0$.  The one exception, obviously, being that
  %% it goes to the very top and stays there.

  %% \begin{answerpart}{Quantum Time (Propagation)}
  %% \end{answerpart}

  %% \begin{answerpart}{How Much Time/Energy Do I Have?}
  %%   We're also going to approach the problem from a different
  %%   uncertainty relation: the time/energy uncertainty.  This one is
  %%   outside the scope of chapter 1...it's just fun...

  %%   \begin{equation}
  %%     E = \frac{1}{2} I \omega^2
  %%   \end{equation}

  %%   Using the same relativity logic as above, we can show that a
  %%   0-energy state (disallowed for other reasons) would be required
  %%   for it to not fall.  Knowing its rotational kinetic energy and a
  %%   time-window in which that (average) energy is true implies the
  %%   ability to compute not only the current $\omega$ but also the
  %%   initial conditions to satisfy the constraints of the fall-time
  %%   equation given the time over which the observation is made.
    
  \end{answerpart}

  
\end{answer}
